\documentclass[11pt]{article}
\usepackage{amsmath, amssymb, amsthm}
\usepackage{algorithm}
\usepackage{algpseudocode}
\usepackage[margin=1in]{geometry}

\title{Rigorous Convergence Analysis of Nesterov's Accelerated Gradient Descent}
\date{}

\begin{document}

\maketitle

\section*{Algorithm Name}
Nesterov's Accelerated Gradient (NAG) Descent for Convex Optimization.

\section*{Mathematical Formulation}
The algorithm maintains two sequences of iterates, $\{x_k\}$ and $\{y_k\}$. The sequence $\{x_k\}$ represents the primary parameter updates, while $\{y_k\}$ represents the extrapolation (momentum) step. 

For $k \ge 1$, the update rules are defined as:
\begin{align*}
    \beta_k &= \frac{k-1}{k+2} \\
    y_k &= x_k + \beta_k (x_k - x_{k-1}) \\
    x_{k+1} &= y_k - \alpha \nabla f(y_k)
\end{align*}
where $\alpha > 0$ is a fixed step size (learning rate) and $\beta_k$ is the deterministic momentum schedule.

\section*{Pseudocode}
\begin{algorithm}[H]
\caption{Nesterov's Accelerated Gradient Descent}
\begin{algorithmic}[1]
\Require Objective function $f$, step size $\alpha \leq 1/L$, initial point $x_1 \in \mathbb{R}^d$, total iterations $T$
\State Initialize $x_0 \gets x_1$ \Comment{Initialization for momentum step}
\For{$k = 1, 2, \dots, T$}
    \State $\beta_k \gets \frac{k-1}{k+2}$
    \State $y_k \gets x_k + \beta_k (x_k - x_{k-1})$ \Comment{Extrapolation step}
    \State $g_k \gets \nabla f(y_k)$ \Comment{Gradient computation at extrapolated point}
    \State $x_{k+1} \gets y_k - \alpha g_k$ \Comment{Gradient descent step}
\EndFor
\Ensure Final iterate $x_{T+1}$
\end{algorithmic}
\end{algorithm}

\section*{Dry Run Example}
Let the objective function be $f(x) = (x-3)^2$. The gradient is $\nabla f(x) = 2(x-3)$. The smoothness constant is $L = 2$.
We choose step size $\alpha = 0.1 \leq 1/L = 0.5$.
Initialize $x_0 = 0$ and $x_1 = 0$.

\textbf{Iteration 1 ($k=1$):}
\begin{itemize}
    \item Momentum: $\beta_1 = \frac{1-1}{1+2} = 0$
    \item Extrapolation: $y_1 = x_1 + 0 \cdot (x_1 - x_0) = 0 + 0 = 0$
    \item Gradient: $\nabla f(y_1) = 2(0-3) = -6$
    \item Update: $x_2 = y_1 - \alpha \nabla f(y_1) = 0 - 0.1(-6) = 0.6$
    \item Objective: $f(x_2) = (0.6 - 3)^2 = 5.76$
\end{itemize}

\textbf{Iteration 2 ($k=2$):}
\begin{itemize}
    \item Momentum: $\beta_2 = \frac{2-1}{2+2} = 0.25$
    \item Extrapolation: $y_2 = x_2 + 0.25(x_2 - x_1) = 0.6 + 0.25(0.6 - 0) = 0.75$
    \item Gradient: $\nabla f(y_2) = 2(0.75-3) = -4.5$
    \item Update: $x_3 = y_2 - \alpha \nabla f(y_2) = 0.75 - 0.1(-4.5) = 1.2$
    \item Objective: $f(x_3) = (1.2 - 3)^2 = 3.24$
\end{itemize}

\textbf{Iteration 3 ($k=3$):}
\begin{itemize}
    \item Momentum: $\beta_3 = \frac{3-1}{3+2} = 0.4$
    \item Extrapolation: $y_3 = x_3 + 0.4(x_3 - x_2) = 1.2 + 0.4(1.2 - 0.6) = 1.44$
    \item Gradient: $\nabla f(y_3) = 2(1.44-3) = -3.12$
    \item Update: $x_4 = y_3 - \alpha \nabla f(y_3) = 1.44 - 0.1(-3.12) = 1.752$
    \item Objective: $f(x_4) = (1.752 - 3)^2 \approx 1.5575$
\end{itemize}

\section*{Assumptions}
The convergence analysis relies on the following standard assumptions:
\begin{enumerate}
    \item \textbf{Convexity:} The objective function $f: \mathbb{R}^d \to \mathbb{R}$ is convex. For all $x, y \in \mathbb{R}^d$:
    $$f(y) \ge f(x) + \langle \nabla f(x), y - x \rangle$$
    \item \textbf{$L$-Smoothness:} The gradient of $f$ is Lipschitz continuous with constant $L > 0$. For all $x, y \in \mathbb{R}^d$:
    $$\|\nabla f(x) - \nabla f(y)\| \le L\|x - y\|$$
    This implies the quadratic upper bound:
    $$f(y) \le f(x) + \langle \nabla f(x), y - x \rangle + \frac{L}{2}\|y - x\|^2$$
    \item \textbf{Existence of Minimum:} The function $f$ achieves its global minimum $f^*$ at some optimal point $x^* \in \mathbb{R}^d$.
    \item \textbf{Valid Step Size:} The learning rate satisfies $\alpha \leq \frac{1}{L}$.
\end{enumerate}

\section*{Main Convergence Theorem}
Under Assumptions 1-4, the sequence $\{x_k\}$ generated by Nesterov's Accelerated Gradient Descent satisfies:
$$f(x_{T+1}) - f(x^*) \le \frac{4}{(T+2)^2} \left( f(x_1) - f(x^*) + \frac{1}{2\alpha} \|x_1 - x^*\|^2 \right)$$
This establishes an $O(1/T^2)$ convergence rate for the objective function value.

\section*{Theoretical Properties}
\begin{itemize}
    \item \textbf{Acceleration Mechanism:} Standard Gradient Descent has a rate of $O(1/T)$ for convex, smooth functions. NAG improves this to $O(1/T^2)$ by evaluating the gradient at an extrapolated point $y_k$ rather than the current iterate $x_k$, effectively utilizing historical gradient information.
    \item \textbf{Optimal Rate:} The $O(1/T^2)$ rate is provably optimal for first-order optimization methods on smooth, convex functions (Nemirovski and Yudin, 1983).
    \item \textbf{Non-monotonicity:} Unlike standard gradient descent, NAG is not strictly a descent method. The objective function value $f(x_k)$ may temporarily increase (exhibiting "ripples" or oscillations) before converging.
\end{itemize}

\section*{Discussion}
The sequence $\beta_k = \frac{k-1}{k+2}$ is specifically engineered to make the error terms telescope perfectly when scaled by a quadratic potential function $t_k^2 \approx k^2/4$. A slight deviation in this schedule can break the acceleration guarantee, highlighting the algorithm's sensitivity to hyperparameter choices.

\section*{Preliminaries (Definitions and Lemmas)}

\textbf{Lemma 1 (Descent Lemma):} For $\alpha \le 1/L$, the update step $x_{k+1} = y_k - \alpha \nabla f(y_k)$ satisfies:
$$f(x_{k+1}) \le f(y_k) - \frac{\alpha}{2} \|\nabla f(y_k)\|^2$$

\textit{Proof of Lemma 1:}
\begin{itemize}
    \item[{[Step 1]}] By the $L$-smoothness assumption, we bound $f(x_{k+1})$ around $y_k$:
    $$f(x_{k+1}) \le f(y_k) + \langle \nabla f(y_k), x_{k+1} - y_k \rangle + \frac{L}{2} \|x_{k+1} - y_k\|^2$$
    \item[{[Step 2]}] Substitute the update rule $x_{k+1} - y_k = -\alpha \nabla f(y_k)$:
    $$f(x_{k+1}) \le f(y_k) - \alpha \|\nabla f(y_k)\|^2 + \frac{L\alpha^2}{2} \|\nabla f(y_k)\|^2$$
    \item[{[Step 3]}] Since $\alpha \le 1/L$, we have $L\alpha^2 \le \alpha$. Thus $\frac{L\alpha^2}{2} \le \frac{\alpha}{2}$. Substituting this yields:
    $$f(x_{k+1}) \le f(y_k) - \frac{\alpha}{2} \|\nabla f(y_k)\|^2$$
\end{itemize}

\textbf{Lemma 2 (One-Step Progress):} For any $z \in \mathbb{R}^d$:
$$f(x_{k+1}) - f(z) \le \frac{1}{2\alpha} \left( \|y_k - z\|^2 - \|x_{k+1} - z\|^2 \right)$$

\textit{Proof of Lemma 2:}
\begin{itemize}
    \item[{[Step 4]}] By the convexity assumption, $f(z) \ge f(y_k) + \langle \nabla f(y_k), z - y_k \rangle$, which rearranges to:
    $$f(y_k) \le f(z) + \langle \nabla f(y_k), y_k - z \rangle$$
    \item[{[Step 5]}] Substitute this upper bound for $f(y_k)$ into Lemma 1:
    $$f(x_{k+1}) \le f(z) + \langle \nabla f(y_k), y_k - z \rangle - \frac{\alpha}{2} \|\nabla f(y_k)\|^2$$
    \item[{[Step 6]}] From the update rule, $\nabla f(y_k) = \frac{1}{\alpha}(y_k - x_{k+1})$. Substitute this into both the inner product and the squared norm:
    $$f(x_{k+1}) - f(z) \le \frac{1}{\alpha} \langle y_k - x_{k+1}, y_k - z \rangle - \frac{1}{2\alpha} \|y_k - x_{k+1}\|^2$$
    \item[{[Step 7]}] Apply the algebraic identity $2\langle a, b \rangle - \|a\|^2 = \|b\|^2 - \|b-a\|^2$ with $a = y_k - x_{k+1}$ and $b = y_k - z$. Note that $b - a = x_{k+1} - z$.
    \item[{[Step 8]}] Factoring out $\frac{1}{2\alpha}$, this directly yields the statement of Lemma 2:
    $$f(x_{k+1}) - f(z) \le \frac{1}{2\alpha} \left( \|y_k - z\|^2 - \|x_{k+1} - z\|^2 \right)$$
\end{itemize}

\section*{Main Proof (Step-by-Step Derivation)}

We construct a potential function proof to establish the accelerated rate.

\begin{itemize}
    \item[{[Step 9]}] Define a sequence $\{t_k\}$ such that $t_k = \frac{k+2}{2}$ for $k \ge 1$, and let $t_0 = 1$. Observe that for $k \ge 1$:
    $$t_k^2 - t_k = \frac{(k+2)^2}{4} - \frac{2k+4}{4} = \frac{k^2+4k+4-2k-4}{4} = \frac{k^2+2k}{4}$$
    Since $t_{k-1}^2 = \frac{(k+1)^2}{4} = \frac{k^2+2k+1}{4}$, we strictly have $t_k^2 - t_k \le t_{k-1}^2$.
    
    \item[{[Step 10]}] Define a reference point $z = \left(1 - \frac{1}{t_k}\right)x_k + \frac{1}{t_k}x^*$. Since $t_k \ge 1$, this is a valid convex combination of $x_k$ and $x^*$.
    
    \item[{[Step 11]}] By the convexity of $f$, applying it to the point $z$ gives:
    $$f(z) \le \left(1 - \frac{1}{t_k}\right)f(x_k) + \frac{1}{t_k}f(x^*)$$
    
    \item[{[Step 12]}] Substitute this convex combination $z$ into Lemma 2:
    $$f(x_{k+1}) - \left[ \left(1 - \frac{1}{t_k}\right)f(x_k) + \frac{1}{t_k}f(x^*) \right] \le \frac{1}{2\alpha} \left( \|y_k - z\|^2 - \|x_{k+1} - z\|^2 \right)$$
    
    \item[{[Step 13]}] Subtract $f(x^*)$ from both sides, and denote the suboptimality gap as $\delta_k = f(x_k) - f(x^*)$. The left side becomes:
    $$\delta_{k+1} - \left(1 - \frac{1}{t_k}\right)\delta_k \le \frac{1}{2\alpha} \left( \|y_k - z\|^2 - \|x_{k+1} - z\|^2 \right)$$
    
    \item[{[Step 14]}] Multiply the entire inequality by $t_k^2$:
    $$t_k^2 \delta_{k+1} - (t_k^2 - t_k)\delta_k \le \frac{t_k^2}{2\alpha} \left( \|y_k - z\|^2 - \|x_{k+1} - z\|^2 \right)$$
    
    \item[{[Step 15]}] Using the property from [Step 9] that $t_k^2 - t_k \le t_{k-1}^2$, and knowing $\delta_k \ge 0$, we can upper bound the left side:
    $$t_k^2 \delta_{k+1} - t_{k-1}^2 \delta_k \le t_k^2 \delta_{k+1} - (t_k^2 - t_k)\delta_k$$
    
    \item[{[Step 16]}] We now analyze the distance terms on the right side. Distribute $t_k$ into the norm for the first term $t_k(y_k - z)$:
    $$t_k y_k - t_k z = t_k y_k - (t_k - 1)x_k - x^*$$
    
    \item[{[Step 17]}] Let us define an auxiliary sequence $u_k = t_{k-1} x_k - (t_{k-1} - 1)x_{k-1}$. We claim that $t_k y_k - (t_k - 1)x_k = u_k$.
    
    \item[{[Step 18]}] From the algorithm definition, $\beta_k = \frac{k-1}{k+2}$. Let us express $\beta_k$ in terms of $t_k$ and $t_{k-1}$:
    $$\frac{t_{k-1} - 1}{t_k} = \frac{\frac{k+1}{2} - 1}{\frac{k+2}{2}} = \frac{\frac{k-1}{2}}{\frac{k+2}{2}} = \frac{k-1}{k+2} = \beta_k$$
    
    \item[{[Step 19]}] Substitute $\beta_k = \frac{t_{k-1} - 1}{t_k}$ into the momentum update $y_k = x_k + \beta_k(x_k - x_{k-1})$, and multiply both sides by $t_k$:
    $$t_k y_k = t_k x_k + (t_{k-1} - 1)(x_k - x_{k-1})$$
    
    \item[{[Step 20]}] Rearrange the terms to isolate $t_k y_k - (t_k - 1)x_k$:
    $$t_k y_k - t_k x_k + x_k = t_{k-1} x_k - (t_{k-1} - 1)x_{k-1}$$
    $$t_k y_k - (t_k - 1)x_k = u_k$$
    This proves the claim from [Step 17].
    
    \item[{[Step 21]}] Substituting [Step 20] into [Step 16], we get the first distance term:
    $$t_k(y_k - z) = u_k - x^*$$
    
    \item[{[Step 22]}] We apply the exact same logic to the second distance term $t_k(x_{k+1} - z)$:
    $$t_k(x_{k+1} - z) = t_k x_{k+1} - (t_k - 1)x_k - x^* = u_{k+1} - x^*$$
    
    \item[{[Step 23]}] Substituting [Step 21] and [Step 22] into the inequality from [Step 14] and [Step 15] yields the fundamental telescoping inequality:
    $$t_k^2 \delta_{k+1} - t_{k-1}^2 \delta_k \le \frac{1}{2\alpha} \left( \|u_k - x^*\|^2 - \|u_{k+1} - x^*\|^2 \right)$$
    
    \item[{[Step 24]}] Sum this inequality over all iterations from $k=1$ to $T$:
    $$\sum_{k=1}^T \left( t_k^2 \delta_{k+1} - t_{k-1}^2 \delta_k \right) \le \frac{1}{2\alpha} \sum_{k=1}^T \left( \|u_k - x^*\|^2 - \|u_{k+1} - x^*\|^2 \right)$$
    
    \item[{[Step 25]}] The sum telescopes perfectly on both sides, leaving only the boundary terms. Furthermore, we drop the non-positive term $-\frac{1}{2\alpha}\|u_{T+1} - x^*\|^2$:
    $$t_T^2 \delta_{T+1} - t_0^2 \delta_1 \le \frac{1}{2\alpha} \|u_1 - x^*\|^2$$
    
    \item[{[Step 26]}] Evaluate the initial conditions. We defined $t_0 = 1$. By the definition in [Step 17], $u_1 = t_0 x_1 - (t_0 - 1)x_0 = 1 \cdot x_1 - 0 \cdot x_0 = x_1$. Substituting these values:
    $$t_T^2 \delta_{T+1} - \delta_1 \le \frac{1}{2\alpha} \|x_1 - x^*\|^2$$
\end{itemize}

\section*{Rate Derivation (With Constants)}
\begin{itemize}
    \item[{[Step 27]}] Rearrange the inequality from [Step 26] to isolate $\delta_{T+1}$:
    $$t_T^2 \delta_{T+1} \le \delta_1 + \frac{1}{2\alpha} \|x_1 - x^*\|^2$$
    
    \item[{[Step 28]}] Substitute the explicit formula for $t_T = \frac{T+2}{2}$, which gives $t_T^2 = \frac{(T+2)^2}{4}$. Substituting $\delta_{T+1} = f(x_{T+1}) - f(x^*)$ and $\delta_1 = f(x_1) - f(x^*)$ yields the final theorem statement:
    $$f(x_{T+1}) - f(x^*) \le \frac{4}{(T+2)^2} \left( f(x_1) - f(x^*) + \frac{1}{2\alpha} \|x_1 - x^*\|^2 \right)$$
\end{itemize}

\section*{Special Cases}
\textbf{Strongly Convex Functions:} 
If the objective function is additionally $\mu$-strongly convex, the standard NAG schedule $\beta_k = \frac{k-1}{k+2}$ is suboptimal. Instead, a constant momentum parameter $\beta = \frac{\sqrt{L} - \sqrt{\mu}}{\sqrt{L} + \sqrt{\mu}}$ is used alongside a step size $\alpha = 1/L$. This adaptation modifies the proof structure to yield a linear convergence rate: $f(x_T) - f(x^*) \le O\left( \left(1 - \sqrt{\frac{\mu}{L}}\right)^T \right)$, exponentially faster than the $O(1/T^2)$ rate for general convex functions.

\textbf{Quadratic Functions:}
For quadratic objectives $f(x) = \frac{1}{2}x^\top Q x - c^\top x$, NAG equates to Chebyshev iteration, and its non-monotonic ripples can be explicitly characterized in terms of the eigenvalues of the Hessian matrix $Q$.

\end{document}